\section{The triple coefficient and arbitrary color products}
\label{sec:coefficient}

We first compute the scalar left by Proposition~\ref{prop:localization},
completing the surface bound. We then convert that bound on isotropic
tests into vanishing for sufficiently long products with arbitrary color
inputs. For the coefficient calculation, put
\[
 t=t_0,\qquad N=t-1+ds_K,\qquad
 A_0=K,\qquad D=dK,\qquad d=3,
 \qquad B=-A_0-D.
\]
The letter \(q\) in the generating series below is a formal variable.
It is unrelated to the realization maps \(q_{g,n}\).

\begin{proposition}\label{prop:scalar}
The bare-triple scalar \(a_{N,t_0}\) in
Proposition~\ref{prop:localization} is nonzero.
Consequently Theorem~\ref{thm:surface} holds.
\end{proposition}

The calculation belongs to the study of tautological Hilbert-scheme
integrals developed by Lehn and Ellingsrud--G\"ottsche--Lehn
\cite[Section~4.3]{Lehn99}, \cite{EGL}. For the symmetric-power and
residue methods used below, compare
\cite[Equation~(9), Theorem~2, and Lemma~3]{MOPCombinatorics}.
The coupling of two Hilbert factors, all marked moments, and the
Euler-shift derivative will be retained explicitly.

We first identify a generating series for this scalar. We then compute
its first nonzero hyperplane coefficient and prove that the resulting
one-variable coefficient cannot vanish.

\subsection{The moving factor}

For a vector bundle \(E\), write \(e_a(E)\) for its Euler class after
adding \(a\) to each Chern root. If \(a\) has invertible constant term,
this convention extends multiplicatively to virtual bundles and
perfect complexes. Inverses are expanded in positive Chow degree.
Let
\[
 M_{r,s}=X^{[r]}\times X^{[s]},
\]
with universal finite subschemes \(Z_1,Z_2\) on \(M_{r,s}\times X\).
Derived Hom expressions in this subsection include pushforward along
the surface \(X\).

On the surviving fixed component, the two kernels are
\[
 I'_1=L_1\mathcal I_{Z_1},\qquad
 I'_2=L_2(-D-x)\mathcal I_{Z_2},
\]
where \(x\) is the hyperplane class of \(|D|\).
Apart from the factor \((1-x)^{\chi(B)}\), the inverse moving Euler
class is
\begin{equation}\label{eq:moving-ratio}
 \begin{aligned}
 R(x)={}&
 \frac{e_1\bigl(((K+A_0)^{[r]})^\vee\bigr)}
 {e_{1-x}\bigl(((K-B)^{[r]})^\vee\bigr)\,
  e_{1-x}(B^{[s]})\,e_{-1+x}((-B)^{[r]})}\\
 &{}\times
 e_{1-x}\bigl(R\Hom(\mathcal O_{Z_1},\mathcal O_{Z_2}(B))\bigr)\,
 e_{1-x}\bigl(R\Hom(\mathcal O_{Z_1},\mathcal O_{Z_2}(K+B))\bigr).
 \end{aligned}
\end{equation}
Here and below only scalar terms of bundles pulled from the original
base can contribute, by Proposition~\ref{prop:localization}.
In particular the omitted inverse Euler factor of
\(R\pi_*(-A_0)\), at weight one, has scalar term one.

To check \eqref{eq:moving-ratio}, the forward moving arrow has
virtual tangent
\[
 R\Hom(\mathcal I_{Z_1},-A_0)
 -
 R\Hom(\mathcal I_{Z_1},\mathcal I_{Z_2}(B)),
\]
at weights \(1\) and \(1-x\), respectively.
Expand each point ideal as \(\mathcal O-\mathcal O_Z\).
The constant term of the second Hom gives
\((1-x)^{\chi(B)}\). Surface duality identifies
\[
 R\Hom(\mathcal O_{Z_1},B)
 \quad\text{with}\quad
 ((K-B)^{[r]})^\vee[-2]
\]
in the convention relevant to the Euler class, and likewise for
\(-A_0\). This gives the first numerator and the first two denominator
factors of \eqref{eq:moving-ratio}.
The reverse moving arrow is
\[
 R\Hom(\mathcal I_{Z_2},\mathcal O_{Z_1}(-B))
\]
at weight \(-1+x\). Its point term gives the final denominator.
Duality changes its cross term to the remaining numerator term.
The cross Hom has virtual rank zero, so changing both its dual and
the sign of its weight introduces no rank sign.

The fixed obstruction Euler classes on this component are
\[
 e\bigl((K^{[r]})^\vee\bigr),
 \qquad
 e_x\bigl(((K-D)^{[s]})^\vee\bigr).
\]
The sign of the second shift is worth recording: the line in the
obstruction bundle is \(K-D-x\) before taking the dual, and hence
has twist \(+x\) afterwards.

\subsection{Marked moments and shifted additivity}

Temporarily replace \(K\) and \(K-D\) in these two fixed Euler
classes by independent divisor variables \(L,H\). Also separate the
Euler shift \(x_e\) from the moving shift \(x_m\). For a surface
class \(z\), let
\[
 \mu_r(z)=
 (\operatorname{pr}_{X^{[r]}})_*
 \bigl([Z_1]\operatorname{pr}_X^*z\bigr),
\]
and define \(\mu_s(z)\) similarly. Consider
\begin{equation}\label{eq:theta-series}
 \Theta(q,p,x_e,x_m)
 =
 \sum_{r,s\geq0}q^{r+s}
 \int_{M_{r,s}}
 e\bigl((L^{[r]})^\vee\bigr)\,
 e_{x_e}\bigl((H^{[s]})^\vee\bigr)\,
 R(x_m)\,
 \exp\bigl(p(\mu_r(z)+\mu_s(z))\bigr).
\end{equation}
The exponential records labeled output insertions with their usual
factorial denominators.

We use the fixed marked recursion of
Proposition~\ref{prop:localization} to regard each finite coefficient
of \eqref{eq:theta-series} as a polynomial in decorated moments
\[
 M_k(\alpha)=\int_X z^k\alpha.
\]
Here \(\alpha\) is a monomial in
\(L,H,A_0,D,K,c_2(T_X)\).
More precisely, carry out the nested-Hilbert recursion while retaining
the marked output slots. Each connected component of the final
equality graph is integrated over one copy of \(X\); if it contains
\(k\) marked slots and decoration \(\alpha\), replace that integration
by \(M_k(\alpha)\). This fixes a polynomial representative before
imposing numerical relations between moments. The recursion and its
projective-bundle signs are those of
\cite[Sections 1 and 3]{EGL}.
The version with arbitrary retained outputs used here is proved in
Lemma~\ref{lem:marked-universality}. This is the interface that keeps
the calculation in relative Chow throughout.

Let \(\mathcal E\) denote the following evaluation:
\begin{equation}\label{eq:bare-moments}
 \mathcal E(M_0(\alpha))=\int_X\alpha,\qquad
 \mathcal E(M_3(1))=1,\qquad
 \mathcal E(M_k(\alpha))=0
 \quad(k>0,\ (k,\alpha)\neq(3,1)).
\end{equation}
Unmarked integrations in this formula are their top-degree numerical
integrals; other degrees give zero.
The scalar normalization agreed with
Proposition~\ref{prop:localization} is
\begin{equation}\label{eq:scalar-coefficient}
 a_{N,t_0}
 =
 [q^N p^{3t}x^t]\,
 (1-x)^{\chi(B)}
 \mathcal E\bigl(\Theta(q,p,x,x)\bigr)
 \big|_{L=K,\ H=K-D}.
\end{equation}
Indeed, \(\mathcal E\) retains precisely the partitions into bare
triples. Passing from the exponential series to labeled outputs
multiplies by the nonzero common factor \((3t)!\).
Every such partition has the same tested value, up to the common
Koszul sign for grouping the two curve axes. Those factors have
already been divided out in Proposition~\ref{prop:localization}.

The next observation allows us to compute \eqref{eq:scalar-coefficient}
using ordinary numerical surface calculations, without assuming that
all formal moments are independent.

\begin{lemma}\label{lem:moment-additivity}
There is a series \(c(q,x_e,x_m)\) such that
\begin{equation}\label{eq:theta-exponential}
 \mathcal E(\Theta)
 =
 \mathcal E(\Theta|_{p=0})
 \exp\bigl(c(q,x_e,x_m)p^3/3!\bigr).
\end{equation}
Moreover, let \(c_H(q)\) be the coefficient of \(M_3(H)\), with the
factor \(p^3/3!\) removed, in the one-moment part of
\(\log\Theta(q,p,0,0)\). Then
\begin{equation}\label{eq:shift-derivative}
 [x_e]c(q,x_e,0)=-c_H(q).
\end{equation}
\end{lemma}

\begin{proof}
Fix a finite coefficient cutoff, so that all polynomials and moment
spaces in this argument are finite-dimensional. On genuine numerical
data, allowing any even inhomogeneous cohomology class \(z\),
\(\Theta\) is multiplicative under disjoint union of surfaces.
The Hilbert schemes split according to the lengths on the components;
all bundles and Hom complexes split; shifted Euler classes are
multiplicative on direct sums; and the incidence exponentials
multiply. These statements hold with both \(x_e\) and \(x_m\)
retained.

Consequently the polynomial for \(\log\Theta\) is additive on
nonnegative integral combinations of genuine moment vectors.
For any finite collection of those vectors, polynomial interpolation
on the nonnegative integer grid extends additivity to their rational
linear span. In particular, its homogeneous components of
moment-count at least two vanish on that span. This is the
disjoint-union argument underlying \cite[Theorem 4.2]{EGL}, applied
coefficientwise to the specified marked polynomial.

The vector in \eqref{eq:bare-moments} belongs to this span.
Start with the genuine unmarked vector obtained by setting \(z=0\).
On a fixed connected surface compare \(z=u+c_0[\mathrm{pt}]\)
with \(z=u\), for rational \(u,c_0\).
The difference changes only the undecorated marked moments, by
\[
 M_k(1)\longmapsto k u^{k-1}c_0.
\]
For every positive-degree decoration the change is zero.
Varying \(u\) and using a finite Vandermonde matrix isolates
\(M_3(1)\) among all the moments within the cutoff.
The logarithm therefore has, at this vector, only an unmarked term
and a term proportional to \(p^3\), proving
\eqref{eq:theta-exponential}.

We next justify differentiating in the Euler shift.
Use the universal polynomial before discarding decoration degrees
larger than the dimension of the surface. In this representative,
the substitution
\[
 H\longmapsto H-x_e
\]
computes the shifted Euler class, since
\[
 e_{x_e}\bigl((H^{[s]})^\vee\bigr)
 =(-1)^s e\bigl((H-x_e)^{[s]}\bigr).
\]
This notation means a base-line twist, not subtraction of a
numerical surface divisor. To realize it, use the constant surface
family over \(\mathbb P^a\), with \(a\) larger than the coefficient
cutoff, and twist \(H\) by \(\mathcal O(-x_e)\).
GRR for the tautological bundle uses \(\exp(H-x_e)\);
the nested-Hilbert recursion and every projection formula commute
with this twist and with moving powers of \(x_e\) out of the
pushforward. It thus gives exactly the substitution in the
 unsimplified moment polynomial. The moving factor is independent of
 the auxiliary \(H\), so its degree cutoff is unaffected.
 More explicitly, an output consisting of \(j\) bare triples requires
 total intrinsic Euler-and-ratio degree \(2(r+s)-2j\) on the Hilbert
 factors. The Euler factors have nonnegative intrinsic degrees.
 Thus moving terms of degree greater than \(2(r+s)\) contribute
 neither to genuine integrals nor to the desired formal coefficient,
 before or after the twist.

It is important to apply additivity to the shifted invariant itself.
The substitution is linear on the moment symbols and preserves their
number in a monomial. For each \(x_e\), the resulting polynomial
still computes the genuine shifted invariant and is additive after
taking its logarithm. Hence every shifted component of moment-count
at least two vanishes on the original genuine span, and in
particular at \(\mathcal E\).
There is no claim that shifting the evaluation vector preserves
that span. For example, an individual term involving \(M_0(H^3)\)
can acquire a nonzero derivative after shifting, but the nonlinear
terms cancel collectively by this additivity.

In the one-moment part of \(p\)-degree three, the derivative followed
by \(\mathcal E\) can only come from
\[
 M_3(H-x_e)=M_3(H)-x_eM_3(1).
\]
Other positive decorations remain positive after a single derivative,
unless the decoration was \(H\) itself. Thus the derivative is
\(-c_H(q)\), proving \eqref{eq:shift-derivative}.
\end{proof}

\subsection{Restriction to two curves}

We first claim that
\begin{equation}\label{eq:zero-unshifted-c}
 c(q,0,x_m)=0.
\end{equation}
Represent very ample \(L,H\) by smooth transverse curves.
At \(x_e=0\), the two Euler classes in
\eqref{eq:theta-series} restrict the integral to
\(L^{(r)}\times H^{(s)}\), with sign \((-1)^{r+s}\).
These are the regular zero loci of the tautological sections:
the zero schemes are the Hilbert schemes of points on the two smooth
curves, of the expected codimensions \(r,s\).
A point-class perturbation of the constant \(z=u\) restricts to zero
on both curves. The Vandermonde test in the preceding proof therefore
annihilates the bare marked coefficient.
The coefficient of the single moment \(M_3(1)\) is a universal scalar,
independent of the numerical line classes, so this proves
\eqref{eq:zero-unshifted-c} identically.

To determine \(c_H\), set both shifts to zero and take
\[
 z=u+\tau Z
\]
for a divisor \(Z\) on the surface. Write \(x',y'\) for the numerical
point divisors on \(L^{(r)},H^{(s)}\).
After the Euler restrictions, the one-body factors of
\eqref{eq:moving-ratio}, including their signs, are
\begin{equation}\label{eq:curve-one-body}
 c_{-1}\bigl((A_0+2D)^{[r]}\bigr)^{-1}
 \,(-1)^s c_{+1}(B^{[s]})^{-1}.
\end{equation}
Here \(c_a(E)=\sum_j a^j c_j(E)\).
For verification, on the ordered cover of a symmetric power,
a tautological line bundle \(M^{[r]}\) has a filtration with roots
\[
 M_i-\delta_i,\qquad
 \delta_i=\sum_{j<i}[\Delta_{ij}].
\]
Two line classes pulled from the same curve position multiply to
zero. Applying this to the low ratio in
\eqref{eq:moving-ratio} reduces its factor at position \(i\) to
\[
 \frac{1}{1+\delta_i-(A_0+2D)_i},
\]
which gives the first term in \eqref{eq:curve-one-body}.

Put \(I=LH\). At each of the \(I\) transverse intersections, the two
universal finite subschemes meet by independently imposing their
point-containment divisors. Their derived Hom has the class of the
structure sheaf of the product of those divisors, twisted by the
determinant normal of the first universal subscheme.
Numerically its class is
\[
 e^{x'}(1-e^{-x'})(1-e^{-y'}).
\]
Its Euler class at weight one is therefore
\[
 \frac{(1+x')(1-y')}{1+x'-y'}.
\]
The two cross Hom terms give the combined factor
\begin{equation}\label{eq:curve-cross}
 \left(\frac{(1+x')(1-y')}{1+x'-y'}\right)^{2I}.
\end{equation}
This calculation also follows by resolving the two transverse
containment divisors and taking their tensor-product resolution;
the twists by surface lines have numerical degree zero at the
fixed intersection point.
The incidence exponential becomes
\begin{equation}\label{eq:curve-incidence}
 \exp\bigl(pu(r+s)+p\tau(ZL\,x'+ZH\,y')\bigr).
\end{equation}

We recall the symmetric-power integral needed to sum these terms.
For a curve of genus \(h\), a line of degree \(d_1\), and its
tautological bundle on the \(r\)-th symmetric power,
\begin{equation}\label{eq:symmetric-chern}
 c_a(\mathrm{line}^{[r]})
 =(1-a x')^{r-d_1-1+h}
 \exp\left(\frac{a\vartheta}{1-a x'}\right),
 \qquad
 \int (x')^{r-j}\frac{\vartheta^j}{j!}=\binom hj.
\end{equation}
Here \(\vartheta\) is the pulled-back Jacobian polarization.
For completeness, the universal divisor sequence and GRR give
\[
 \ch(\mathrm{line}^{[r]})
 =d_1+1-h+
   (r-d_1-1+h+\vartheta)\exp(-x').
\]
Indeed the universal divisor class has terms
\(r[\mathrm{pt}]+x'+\ell\), where the mixed class satisfies
\(\int_{\mathrm{curve}}\ell^2/2=-\vartheta\).
Converting this Chern character to the total Chern class gives the
first formula in \eqref{eq:symmetric-chern}.
For the second, pull to the ordered cover and expand the polarization
in a symplectic basis. Intersecting the point divisors first leaves
\(j\) factors; choosing \(j\) distinct symplectic pairs contributes
one after the factorials, giving \(\binom hj\).

Let
\[
 d_1=(A_0+2D)L,\qquad d_2=BH,\qquad
 2h_1-2=L(L+K),\qquad 2h_2-2=H(H+K).
\]
By \eqref{eq:symmetric-chern}, low-curve integration of the first
factor in \eqref{eq:curve-one-body}, times any series in \(x'\),
is coefficient extraction against
\begin{equation}\label{eq:curve-integration}
 [ (x')^r]\,
 (1+x')^{d_1+1-2h_1-r}(1+2x')^{h_1}.
\end{equation}
High-curve integration has the analogous expression after replacing
\(y'\) by \(-y'\), which absorbs the sign \((-1)^s\).
Accordingly \eqref{eq:curve-cross} becomes
\[
 \left(\frac{(1+x')(1+y')}{1+x'+y'}\right)^{2I},
\]
and the nonconstant part of \eqref{eq:curve-incidence} becomes
\(\exp(p\tau(ZL\,x'-ZH\,y'))\).

The elementary Lagrange identity used for each variable is
\begin{equation}\label{eq:lagrange-hilbert}
 \sum_{r\geq0}Q^r [x^r]\frac{F(x)}{(1+x)^r}
 =F(w)\frac{1+w}{1+2w},
 \qquad w(1+w)=Q,\quad w(0)=0.
\end{equation}
One proof is to express coefficient extraction as the formal residue
at zero: summing the geometric series gives
\(F(x)(1+x)/(x(1+x)-Q)\), whose small root is \(w\);
its residue there is the right side.
The identity holds with other formal variables in the coefficients,
so it can be applied successively to \(x'\) and \(y'\).
Using \(Q=q\exp(pu)\), both small roots equal \(w\).
Equations \eqref{eq:curve-integration}--\eqref{eq:lagrange-hilbert}
therefore give
\begin{align}
 \log\Theta(q,p,0,0)
 ={}&
 (d_1+d_2+4-2h_1-2h_2+4I)\log(1+w) \notag\\
 &+(h_1+h_2-2-2I)\log(1+2w)
   +p\tau(ZL-ZH)w .
 \label{eq:theta-curves}
\end{align}
This equation is the genuine numerical evaluation with
\(z=u+\tau Z\). It determines the universal coefficients used below.

\subsection{The first hyperplane coefficient}

Let \(w_0(1+w_0)=q\), with \(w_0(0)=0\), and let
\(\mathcal D=q\partial_q\).
By Lemma~\ref{lem:moment-additivity}, the numerical evaluation of
\(\log\Theta\) on a genuine moment vector equals that of its
one-moment part: every component with at least two moment factors
vanishes there. Its term of \(p\)-degree three is a universal scalar
combination of the moments \(M_3(\alpha)\). The term involving
\(M_3(H)\) is \(c_H(q)M_3(H)p^3/3!\).
Since
\[
 [u^2\tau]M_3(H)=3ZH,
\]
we can identify \(c_H\) by varying the very ample \(H\), with the
other divisor classes fixed. Only degree-one decorations can
contribute to \(u^2\tau\), and \(H\) is the only varying such
decoration. For example, one may take arbitrarily large integral
multiples of a fixed ample line for \(H\), with \(Z\) fixed ample.
The first two terms of \eqref{eq:theta-curves} contain no \(\tau\).
The coefficient of \(p^3u^2\tau(ZH)\) in its last term is
\(-\mathcal D^2w_0/2\). Comparing factorials yields
\[
 c_H(q)=-\mathcal D^2w_0.
\]
Together with \eqref{eq:shift-derivative} and
\eqref{eq:zero-unshifted-c}, this proves
\begin{equation}\label{eq:first-c-coefficient}
 [x]c(q,x,x)
 =\mathcal D^2w_0
 =\frac{q(1+2q)}{(1+4q)^{3/2}}.
\end{equation}
There is no contribution from differentiating \(x_m\), because
\(c(q,0,x_m)\) is identically zero.

The unmarked calculation is equally explicit.
Set \(p=0\) in \eqref{eq:theta-curves}.
Every finite coefficient is polynomial in the intersection numbers
of \(L,H,A_0,D,K\), so the formula continues from very ample
\(L,H\) to \(L=K,\ H=K-D\).
This can also be seen by varying sufficiently large multiples of
\(K\) in both variables and then using polynomial identity.
The adjunction expressions for \(h_1,h_2\) are polynomial
expressions in this continuation, not assertions that a negative
line has a smooth effective representative.
Substitution gives
\begin{equation}\label{eq:unmarked-theta}
 \mathcal E(\Theta(q,0,0,0))
 \big|_{L=K,\ H=K-D}
 =
 (1+w_0)^{ds_K}
 (1+2w_0)^{d(d+1)s_K/2}.
\end{equation}

\subsection{A residue with a definite sign}

We now prove Proposition~\ref{prop:scalar}.
By \eqref{eq:theta-exponential}, the \(p^{3t}\)-coefficient is
\(c(q,x,x)^t/(6^t t!)\) times the unmarked series.
The series \(c(q,x,x)\) has zero constant term in \(x\).
Thus the coefficient of \(x^t\) uses its linear term in every factor,
and uses only the constant terms in \(x\) of the unmarked series and
of \((1-x)^{\chi(B)}\).
Equations \eqref{eq:scalar-coefficient},
\eqref{eq:first-c-coefficient}, and \eqref{eq:unmarked-theta} give
\begin{equation}\label{eq:scalar-final-series}
 a_{N,t_0}
 =
 \frac{1}{6^t t!}
 [q^{ds_K-1}]
 (1+w_0)^{ds_K}(1+2q)^t(1+4q)^E,
 \qquad
 E=\frac{d(d+1)s_K}{4}-\frac{3t}{2}.
\end{equation}

Put \(v=1+2w_0\), so that \(q=(v^2-1)/4\).
The coefficient residue in \eqref{eq:scalar-final-series} becomes
\begin{equation}\label{eq:scalar-residue}
 2^{ds_K-t-1}
 \operatorname*{Res}_{v=1}
 \frac{(1+v^2)^t v^{2E+1}}{(v-1)^{ds_K}}\,dv.
\end{equation}
For the present parameters,
\[
 ds_K=24\chi,\qquad
 t=25\chi-1,\qquad
 E=\frac{3-27\chi}{2},\qquad
 \chi=(b-1)^2\geq16.
\]
Hence the rational differential in \eqref{eq:scalar-residue} has
finite poles only at zero and one.
Its rational coefficient has degree \(2-\chi\leq-14\) at infinity,
so its residue at infinity is zero.
The residue at zero is
\[
 [v^{27\chi-5}]
 (1+v^2)^{25\chi-1}(1-v)^{-24\chi}.
\]
This is strictly positive: the second factor has positive
coefficients in every nonnegative degree, the first has nonnegative
coefficients and constant term one, and \(27\chi-5\geq0\).
The residue at one is its negative and therefore is nonzero.
Both scalar prefactors in
\eqref{eq:scalar-final-series}--\eqref{eq:scalar-residue} are nonzero.
This proves \(a_{N,t_0}\neq0\).

Finally, Proposition~\ref{prop:localization} gives
\(a_{N,t_0}T_\gamma^{t_0}=0\) for the isotropic cubic contraction
\(T_\gamma\) appearing in Theorem~\ref{thm:surface}.
The scalar is a nonzero rational number and is invertible in the
coefficient algebra. Thus \(T_\gamma^{t_0}=0\), which completes
the proof of that theorem.

\subsection{From isotropic tests to arbitrary color products}

The surface theorem now gives nilpotence for every isotropic tensor
test. To use it in the product defining the
counterexample, we need a vanishing statement for arbitrary color
inputs. We return to a smooth genus-$b$ curve family, choosing a rational
divisor of relative degree one and its pure projector as in
Proposition~\ref{prop:tensor-test}. We use the operation algebra $\mathcal A_k(E)$ of
Section~\ref{sec:tensors} for an ordinary rational color space $E$.
A \emph{pure cubic entry} is the class of the projected small diagonal
in $\mathcal A_3(E)$ with a specified triple of color inputs; products
use disjoint curve slots.

We will fix an arbitrary color map $E\to W_b$, apply the surface theorem
to a subspace with isotropic image, and use the following exterior-algebra
lemma to annihilate every sufficiently long product. The final
application of faithful polarization will then give a Chow identity
with the same base-support bound.

\begin{lemma}[An exterior ideal]
\label{lem:exterior-ideal}
Let \(F\subset H\) be vector spaces of dimensions \(2r\) and \(D\), and
let \(\Omega\in\bigwedge^2F\) be nondegenerate.  If \(0\le t\le r\), then
the ideal generated by \(\Omega^t\) in \(\bigwedge H\) contains
\(\bigwedge^qH\) for every \(q\ge D-r+t\).
\end{lemma}

\begin{proof}
The linear symplectic Lefschetz decomposition gives surjectivity of
\[
 \Omega^t\wedge-:
 \bigwedge^{k-2t}F\longrightarrow\bigwedge^kF
 \qquad(k\ge r+t).
\]
Equivalently, decompose the exterior algebra into primitive Lefschetz
strings; every term in those degrees is at least \(t\) positions up its
string.  Split \(H=F\oplus F'\).  A nonzero summand of
\(\bigwedge^qH\) has \(F'\)-degree at most \(D-2r\), hence \(F\)-degree
at least \(q-(D-2r)\ge r+t\).  It therefore lies in the stated ideal.
\end{proof}

\begin{lemma}[Arbitrary color products]
\label{lem:large-vertex-test}
Let a smooth curve family of genus \(b\ge5\), with a section, be over a
smooth irreducible quasi-projective base \(S\) satisfying
\(\dim S<(b-1)^2-1\).  Put
\[
 t_b=25(b-1)^2-1.
\]
Let \(E\) be an \(m\)-dimensional color space, and let \(e\) be an even
integer with \(e\ge m/2-2\) and \(e\le m/2\).
Write \(D_E=\binom m3\) and \(d_e=\binom e3\).
If \(t_b\le d_e/2\), every product of
\[
 q\ge D_E-d_e/2+t_b
\]
pure cubic entries whose colors lie in \(E\) vanishes modulo support in
base codimension at least \(b-1\).
\end{lemma}

\begin{proof}
Fix a color test \(z:E\to W_b\).  The pullback to \(E\) of the symplectic
form has a totally isotropic subspace of dimension at least \(m/2\):
its radical together with a Lagrangian in its symplectic quotient has
that dimension.  Choose an \(e\)-dimensional subspace \(E'\) of it.
Its image under \(z\) is isotropic.

Choose an invertible skew matrix \(\gamma\) in a basis of \(E'\).
Grouping the ordered triples in Theorem~\ref{thm:surface}, the coefficients
of the resulting quadratic expression in cubic generators are a
nonzero common scalar times the three-by-three minors of \(\gamma\).
Thus this expression is a nondegenerate bivector
\[
 \Omega\in\bigwedge^2\bigl(\bigwedge^3E'\bigr)
\]
up to a nonzero scalar.  Indeed, the third compound matrix is invertible,
and its transpose is its negative because three is odd.
In particular \(d_e\) is even.  Theorem~\ref{thm:surface} says that the
map from the exterior algebra on \(\bigwedge^3 E\) to the odd cubic
coefficients sends \(\Omega^{t_b}\) to zero in the specified support
quotient.

Lemma~\ref{lem:exterior-ideal}, with
\(F=\bigwedge^3E'\) and \(H=\bigwedge^3E\), now kills every element of
the asserted degree.  This holds for every \(z\); although \(E'\) can
depend on \(z\), the degree threshold does not.
Proposition~\ref{prop:tensor-test} therefore proves the Chow statement.
Its finite coefficient test also ensures that the resulting support
still has codimension at least \(b-1\).
\end{proof}


The smooth-family argument now gives vanishing of long products of
pure cubic entries for arbitrary color inputs, modulo the stated
base-support bound. The remaining task is
to compare the compact candidate with these pure entries and control
the errors along the boundary.
